\documentclass[a4paper,10pt,twocolumn]{ltjarticle}
\usepackage{kaitoushu}

\pgfplotsset{
  myaxis/.style={
    axis lines=middle,
    xlabel={$x$}, ylabel={$y$},
    every axis x label/.style={at={(ticklabel* cs:1)},anchor=west},
    every axis y label/.style={at={(ticklabel* cs:1)},anchor=south},
    tick label style={font=\scriptsize},
    label style={font=\small},
    width=5.5cm, height=5cm,
    grid=none,
    samples=200,
  }
}

\begin{document}

\problemrange{5章：三角関数 — Basic (p.69) 加法定理・倍角・合成（問318--325）}



\mondai{318}{次の値を求めよ．}
\kotae{
加法定理：$\sin(\alpha\pm\beta)=\sin\alpha\cos\beta\pm\cos\alpha\sin\beta$，$\cos(\alpha\pm\beta)=\cos\alpha\cos\beta\mp\sin\alpha\sin\beta$

(1) $\sin75°=\sin(45°+30°)=\dfrac{\sqrt{2}}{2}\cdot\dfrac{\sqrt{3}}{2}+\dfrac{\sqrt{2}}{2}\cdot\dfrac{1}{2}=\dfrac{\sqrt{6}+\sqrt{2}}{4}$

(2) $\cos15°=\cos(45°-30°)=\dfrac{\sqrt{2}}{2}\cdot\dfrac{\sqrt{3}}{2}+\dfrac{\sqrt{2}}{2}\cdot\dfrac{1}{2}=\dfrac{\sqrt{6}+\sqrt{2}}{4}$

(3) $\tan\dfrac{5\pi}{12}=\tan(75°)=\tan(45°+30°)=\dfrac{1+\frac{1}{\sqrt{3}}}{1-\frac{1}{\sqrt{3}}}=\dfrac{\sqrt{3}+1}{\sqrt{3}-1}=2+\sqrt{3}$
}

\mondai{319}{$\alpha$が第2象限で$\sin\alpha=\dfrac{4}{5}$，$\beta$が第1象限で$\cos\beta=\dfrac{5}{13}$のとき，$\sin(\alpha+\beta)$，$\cos(\alpha-\beta)$ を求めよ．}
\kotae{
$\cos\alpha=-\dfrac{3}{5}$（第2象限），$\sin\beta=\dfrac{12}{13}$（第1象限）

$\sin(\alpha+\beta)=\dfrac{4}{5}\cdot\dfrac{5}{13}+\left(-\dfrac{3}{5}\right)\cdot\dfrac{12}{13}=\dfrac{20}{65}-\dfrac{36}{65}=-\dfrac{16}{65}$

$\cos(\alpha-\beta)=\left(-\dfrac{3}{5}\right)\cdot\dfrac{5}{13}+\dfrac{4}{5}\cdot\dfrac{12}{13}=-\dfrac{15}{65}+\dfrac{48}{65}=\dfrac{33}{65}$
}

\mondai{320}{$\sin\alpha=\dfrac{1}{\sqrt{5}}$，$\sin\beta=\dfrac{1}{\sqrt{10}}$（$\alpha,\beta$は鋭角）のとき，$\alpha+\beta$を求めよ．}
\kotae{
$\cos\alpha=\dfrac{2}{\sqrt{5}}$，$\cos\beta=\dfrac{3}{\sqrt{10}}$

$\sin(\alpha+\beta)=\dfrac{1}{\sqrt{5}}\cdot\dfrac{3}{\sqrt{10}}+\dfrac{2}{\sqrt{5}}\cdot\dfrac{1}{\sqrt{10}}=\dfrac{3}{\sqrt{50}}+\dfrac{2}{\sqrt{50}}=\dfrac{5}{5\sqrt{2}}=\dfrac{1}{\sqrt{2}}$

$\alpha+\beta$ は鋭角の和なので $0<\alpha+\beta<\pi$．$\sin(\alpha+\beta)=\dfrac{\sqrt{2}}{2}$ より

$\alpha+\beta=\dfrac{\pi}{4}$ または $\dfrac{3\pi}{4}$

$\cos(\alpha+\beta)=\dfrac{2}{\sqrt{5}}\cdot\dfrac{3}{\sqrt{10}}-\dfrac{1}{\sqrt{5}}\cdot\dfrac{1}{\sqrt{10}}=\dfrac{5}{\sqrt{50}}=\dfrac{1}{\sqrt{2}}>0$ より

$\alpha+\beta=\dfrac{\pi}{4}$
}

\mondai{321}{2倍角の公式を用いて，次の値を求めよ．}
\kotae{
$\sin2\alpha=2\sin\alpha\cos\alpha$，$\cos2\alpha=1-2\sin^{2}\alpha=2\cos^{2}\alpha-1$

(1) $\sin\dfrac{\pi}{8}$

$\cos\dfrac{\pi}{4}=1-2\sin^{2}\dfrac{\pi}{8}$，$\dfrac{\sqrt{2}}{2}=1-2\sin^{2}\dfrac{\pi}{8}$

$\sin^{2}\dfrac{\pi}{8}=\dfrac{2-\sqrt{2}}{4}$，$\sin\dfrac{\pi}{8}=\dfrac{\sqrt{2-\sqrt{2}}}{2}$

(2) $\cos\dfrac{\pi}{8}$

$\cos\dfrac{\pi}{4}=2\cos^{2}\dfrac{\pi}{8}-1$，$\cos^{2}\dfrac{\pi}{8}=\dfrac{2+\sqrt{2}}{4}$

$\cos\dfrac{\pi}{8}=\dfrac{\sqrt{2+\sqrt{2}}}{2}$
}

\mondai{322}{$\theta$が第3象限で$\cos\theta=-\dfrac{3}{4}$のとき，$\sin2\theta$，$\cos2\theta$，$\tan2\theta$を求めよ．}
\kotae{
$\sin^{2}\theta=1-\dfrac{9}{16}=\dfrac{7}{16}$．第3象限で$\sin<0$：$\sin\theta=-\dfrac{\sqrt{7}}{4}$

$\sin2\theta=2\cdot\left(-\dfrac{\sqrt{7}}{4}\right)\!\left(-\dfrac{3}{4}\right)=\dfrac{3\sqrt{7}}{8}$

$\cos2\theta=2\cos^{2}\theta-1=2\cdot\dfrac{9}{16}-1=\dfrac{1}{8}$

$\tan2\theta=\dfrac{\sin2\theta}{\cos2\theta}=\dfrac{3\sqrt{7}/8}{1/8}=3\sqrt{7}$
}

\mondai{323}{三角関数の合成を行え．}
\kotae{
$a\sin\theta+b\cos\theta=\sqrt{a^{2}+b^{2}}\sin(\theta+\varphi)$（$\tan\varphi=\dfrac{b}{a}$）

(1) $\sin\theta+\cos\theta=\sqrt{2}\sin\!\left(\theta+\dfrac{\pi}{4}\right)$

(2) $\sin\theta-\cos\theta=\sqrt{2}\sin\!\left(\theta-\dfrac{\pi}{4}\right)$

(3) $\sqrt{3}\sin\theta+\cos\theta=2\sin\!\left(\theta+\dfrac{\pi}{6}\right)$

(4) $\sin\theta-\sqrt{3}\cos\theta=2\sin\!\left(\theta-\dfrac{\pi}{3}\right)$
}

\mondai{324}{$0\leqq\theta<2\pi$ のとき，$y=\sin\theta+\cos\theta$ の最大値と最小値を求めよ．}
\kotae{
$y=\sqrt{2}\sin\!\left(\theta+\dfrac{\pi}{4}\right)$

$-1\leqq\sin\!\left(\theta+\dfrac{\pi}{4}\right)\leqq1$ より

最大値$\sqrt{2}$（$\theta=\dfrac{\pi}{4}$），最小値$-\sqrt{2}$（$\theta=\dfrac{5\pi}{4}$）
}

\mondai{325}{$0\leqq\theta<2\pi$ のとき，$\sqrt{3}\sin\theta-\cos\theta=1$ を解け．}
\kotae{
$\sqrt{3}\sin\theta-\cos\theta=2\sin\!\left(\theta-\dfrac{\pi}{6}\right)=1$

$\sin\!\left(\theta-\dfrac{\pi}{6}\right)=\dfrac{1}{2}$

$\theta-\dfrac{\pi}{6}=\dfrac{\pi}{6},\;\dfrac{5\pi}{6}$

$\theta=\dfrac{\pi}{3},\;\pi$
}

\end{document}
